% Current method chapter, 06 October 2026.
% Dependencies: the existing main.tex preamble (amsmath, TikZ, natbib, hyperref).
% Source/evidence are read-only. This chapter distinguishes specification,
% executable simulation, and empirical confirmation.
\chapter{Phương pháp dựa trên Geo-I và truy hồi POI tại thiết bị}
\label{ch:method}
\label{ch:current-method}

\section{Phạm vi và cấu hình hiện tại}
\label{sec:current-method-scope}

Phương pháp chạy ở thiết bị: bảo vệ những lần đọc vị trí bằng Geo-I, dùng lịch sử đã bảo vệ để chọn các tọa độ truy vấn khả thi trên đường, rồi phục hồi câu trả lời bằng xếp hạng cục bộ. Geo-I là nền tảng riêng tư; REM là cơ chế lấy mẫu điểm tham chiếu trên một miền đường công khai cố định. Tập truy vấn giả và bước lấy POI giúp giữ chất lượng dịch vụ, nhưng không tự tạo thêm một bảo đảm riêng tư mới.

Cấu hình được xác nhận trong vòng mới mang tên \textbf{Geo-I / REM, $L=30$}. Nó giữ bộ chọn \texttt{legacy\_l10}: chữ ``10'' chỉ độ sâu bảng POI mà bộ chọn dùng, còn $L=30$ là số POI tối đa máy chủ trả \emph{cho mỗi loại tại mỗi tọa độ truy vấn}. Vẫn gửi $K=5$ tọa độ và trả tối đa $k=5$ POI tại thiết bị. $L$ không phải các mã cấu hình 30/67 trong tài liệu lịch sử. Hai vòng thay mục tiêu bộ chọn Q đã không qua điều kiện chọn; cấu hình được giữ chỉ tăng độ sâu phản hồi, không thay cơ chế Geo-I hay Q đã sinh. Các kết quả thực nghiệm được trình bày riêng, thay vì suy chất lượng bảo vệ từ tên cấu hình.\footnote{Cấu hình đã chốt: \path{artifacts/benchmarks/qplanner_response_depth_generalization_20261006_v1/recommended_configuration.json}. Lịch sử giữ cả thử nghiệm thất bại: \path{docs/reviews/2026-10-06_qplanner_iteration_log.md}.}

\section{Luồng xử lý và ranh giới quan sát}
\label{sec:current-method-flow}

Gọi $\mathcal C$ là ngữ cảnh công khai: phép chiếu tọa độ, mạng đường có hướng, miền trạng thái, danh mục và phiên bản POI, prior, tham số, lịch gửi và chính sách cấp phiên. Tại sự kiện công khai $t$, thiết bị có vị trí thật $x_t$ và nhu cầu riêng $u_t$ gồm mục đích, loại POI, bán kính hoặc đích đến. Hai đầu vào này có đường xử lý khác nhau. $x_t$ chỉ đi vào cơ chế bảo vệ khi lịch đọc và bộ lọc ngân sách cho phép; $u_t$ chỉ đi vào xếp hạng sau khi nhận phản hồi.

\begin{figure}[htbp]
\centering
\resizebox{\textwidth}{!}{%
\begin{tikzpicture}[
 font=\small,
 b/.style={draw=reportblue,rounded corners,align=center,text width=37mm,minimum height=13mm,fill=white},
 input/.style={b,font=\footnotesize,text width=29mm,minimum height=13mm},
 a/.style={-{Latex[length=2mm]},thick,draw=reportblue},
 local/.style={a,draw=reportgreen},
 pub/.style={a,dashed,draw=black!55}]
\filldraw[rounded corners,fill=green!3,draw=reportgreen] (-.3,.95) rectangle (10.45,-10.05);
\filldraw[rounded corners,fill=orange!8,draw=orange!70!black,line width=1pt] (-.05,-1.4) rectangle (10.2,-7.6);
\filldraw[rounded corners,fill=red!4,draw=reportred] (10.9,.95) rectangle (15.15,-10.05);
\node[font=\small\bfseries,text=reportgreen] at (5.0,.6) {THIẾT BỊ: VÙNG TIN CẬY};
\node[font=\small\bfseries,text=reportred,align=center] at (13.0,.45) {MÁY CHỦ\\VÙNG QUAN SÁT};
\node[font=\small\bfseries,text=orange!70!black,align=center,text width=43mm] at (7.7,-2.35) {PHƯƠNG PHÁP\\CỦA LUẬN VĂN};
\node[input,fill=green!8] (gps) at (1.7,-.35) {Đầu vào riêng:\\GPS thật $x_t$};
\node[input,fill=reportgray] (public) at (5.2,-.35) {Đầu vào công khai:\\bản đồ, POI, prior, lịch};
\node[input,fill=green!8] (privatequery) at (8.6,-.35) {Đầu vào riêng $u_t$:\\mục đích, loại,\\bán kính, đích};
\node[b,fill=orange!15] (filter) at (2.2,-2.95) {1. Lịch và cap\\2. REM / tái dùng neo $Z$};
\node[b,fill=reportlight] (belief) at (2.2,-4.75) {3. Ước lượng từ lịch sử\\$Z$ và nhánh đã bảo vệ};
\node[b,fill=reportlight] (planner) at (2.2,-6.65) {4. Chọn $K=5$ Q\\theo đường có thể tới};
\node[b,fill=reportlight] (rank) at (7.85,-6.65) {6. Hợp POI đã nhận\\xếp hạng bằng $x_t,u_t$};
\node[b,draw=reportred,fill=red!8] (server) at (13.0,-6.65) {5. Lấy ứng viên mọi loại\\$L=30$ mỗi loại / mỗi Q};
\node[b,fill=green!8] (answer) at (7.85,-9.05) {Đầu ra cục bộ:\\tối đa $k=5$ POI phù hợp};
\node[align=center,font=\scriptsize,text width=43mm] at (2.2,-9.05) {Bản tin: tọa độ Q, thời điểm,\\mọi loại POI và $L$ cố định.\\Không gửi $Z$ hay nhu cầu thật.};
\node[align=center,font=\scriptsize,text width=38mm] at (13.0,-2.75) {Có thể lưu tọa độ, thứ tự,\\thời điểm và kích thước\\request / response.};
\draw[a] (gps)--(filter);
\draw[a] (filter)--(belief);
\draw[a] (belief)--(planner);
\draw[a] (planner.south)--(2.2,-7.95)--node[pos=.33,above,font=\scriptsize,fill=green!3]{Q gửi ra}(13.0,-7.95)--(server.south);
\draw[a] (server.west)--(rank.east);
\draw[local] (privatequery.south)--(8.6,-1.15)--(10.0,-1.15)--(10.0,-5.55)--(rank.north);
\draw[local] (gps.east)--(3.55,-.35)--(3.55,-1.15)--(4.65,-1.15)--(4.65,-5.55)--(rank.north west);
\node[font=\scriptsize,fill=orange!8,rotate=90] at (4.65,-3.5) {GPS dùng local};
\draw[local,preaction={draw=white,line width=3pt}] (rank)--(answer);
\draw[pub] (public.south)--(5.2,-4.75)--(belief.east);
\draw[pub] (5.2,-4.75)--(5.2,-5.65)--(planner.north east);
\end{tikzpicture}}
\caption{Luồng hiện tại. Vùng cam là phương pháp tại thiết bị. Điểm tham chiếu $Z$ phục vụ nội bộ; Q là tọa độ gửi ra, còn POI là kết quả dịch vụ. Nhu cầu thật không chọn Q hoặc kích hoạt truy vấn mạng.}
\label{fig:current-method-flow}
\end{figure}

Ba ký hiệu cần phân biệt: \textbf{$Z_t$} là điểm tham chiếu đã bảo vệ, có thể giữ lại qua nhiều lần gửi; \textbf{$Q_t=(q_{t,1},\ldots,q_{t,K})$} là các tọa độ truy vấn từ trạng thái đã bảo vệ; \textbf{$P_t$} là các POI do dịch vụ trả về. Q không phải năm mẫu REM độc lập quanh GPS, không buộc chứa GPS thật và cũng không buộc mọi tọa độ khác nhau. Cơ chế không cam kết GPS thật không bao giờ trùng ngẫu nhiên với một Q.

Đối thủ có thể thấy tọa độ Q, thứ tự trong bản tin, các mốc công khai, loại POI công khai, $L$, phản hồi và kích thước. Quan sát còn có thể chứa giờ mở/đóng phiên theo giao thức đánh giá. GPS, $Z$, nhánh giữ/tạo mới, số dư, khóa ngẫu nhiên, $u_t$ và câu trả lời sau xếp hạng nằm ngoài quan sát này. Account/IP, click, đặt chỗ và điều hướng tiếp theo không được ẩn bởi cơ chế tọa độ. Nếu ứng dụng công bố các thông tin đó hoặc tải bản đồ theo GPS thật, phải bổ sung chúng vào mô hình đe dọa.

Adapter có thể xáo thứ tự Q và bỏ ID cố định, nhưng vẫn có thể nối Q bằng hình học. Vòng xác nhận L30 giữ quan sát slot có thứ tự và cả đối thủ nối bằng Hungarian; không quy kết kết quả đó cho khả năng ``không thể theo dõi'' của việc xáo thứ tự.\footnote{Ranh giới API: \path{benchmark/geoi_lbs.py}, \path{benchmark/budgeted_geoi_lbs.py}, \path{benchmark/query_purpose.py}, \path{benchmark/query_order.py}. Hợp đồng vòng L30: \path{experiments/qplanner_study_20261006_v2.py}, hàm \texttt{run\_family}.}

\section{REM và phép thử tái sử dụng có nhiễu}
\label{sec:current-rem}

Miền $V$ chứa các trạng thái làn--độ tiến được dựng từ bản đồ công khai trước khi xử lý chuyến. Một trạng thái $v$ có tọa độ $g(v)$. GPS được đưa qua phép chiếu cố định $\Pi$; cơ chế dùng $x=\Pi(\mathrm{GPS})$ liên tục, không làm tròn GPS thành nút đường trước khi lấy mẫu. Khoảng cách riêng tư là $d_E(x,x')=\|x-x'\|_2$ tính bằng mét trong hệ phẳng này. Khoảng cách đường có hướng dùng cho dịch vụ là đại lượng khác, có thể bất đối xứng hoặc vô hạn.

Với $u>0$ theo đơn vị $\mathrm m^{-1}$, kernel REM lý tưởng là
\begin{equation}
 R_u(v\mid x)=\frac{\exp[-u\,d_E(x,g(v))/2]}
 {\sum_{w\in V}\exp[-u\,d_E(x,g(w))/2]},\qquad v\in V.
 \label{eq:current-rem}
\end{equation}
Toàn bộ $V$ giữ cố định, không cắt bằng bán kính quanh GPS thật. REM sử dụng điểm số Euclid trên miền đường; không đồng nhất nó với GEM dùng khoảng cách đồ thị của Takagi và cộng sự.\footnote{S. Takagi và cộng sự, \emph{Geo-Graph-Indistinguishability}, bản đầy đủ: \url{https://arxiv.org/abs/2010.13449}. Mã REM và giới hạn số học: \path{core/mechanisms.py}.} Geo-I và các tính chất hậu xử lý là nền tảng kế thừa~\cite{andres2013geoi}.

Lần đọc đầu tạo $Z$ bằng REM. Những lần đọc sau lấy $h=Z$ cuối cùng làm tham chiếu, không lấy GPS cũ làm tham chiếu, rồi thử
\begin{equation}
 d_E(x_t,h)+\eta_t\leq\theta,
 \qquad \eta_t\sim\operatorname{Lap}(1/u).
 \label{eq:current-reuse}
\end{equation}
Nếu đạt, giữ $Z_t=h$; nếu không, lấy $Z_t$ mới bằng (\ref{eq:current-rem}). Ngưỡng $\theta=200$ m là công khai. Quyết định không phải ``ra ngoài200 m thì chắc chắn tạo mới'': nhiễu làm cả hai nhánh có xác suất. Tái sử dụng vẫn tốn ngân sách phép thử; đây là cách khởi tạo cụ thể của cơ chế dự đoán có kiểm tra riêng tư~\cite{chatzikokolakis2014predictive}, không phải primitive mới. Nó hạn chế số mẫu mới trong một chân trời hữu hạn, không ngăn tấn công trung bình vô hạn hay ghi nhớ bền vững mỗi lần quay về nhà.

\section{Bộ lọc ngân sách qua các phiên liên kết}
\label{sec:current-budget}

Chính sách công khai dành trước $N$ slot trong một epoch cho cùng chủ thể được quản lý bằng một sổ. $C$ là cap hiệu lực của cả epoch; $H$ điều khiển cách chia chi phí trong engine. Đặt
\begin{equation}
 U=2H-1,\qquad u=\frac{C}{N U},\qquad
 C_{\mathrm{session}}=Uu=\frac{C}{N},\qquad
 B_{\mathrm{session}}=2Hu.
 \label{eq:current-budget}
\end{equation}
$B_{\mathrm{session}}$ là tham số danh nghĩa mà engine nhận; $C_{\mathrm{session}}$ là cap hiệu lực. Không hoán đổi hai con số này khi so sánh. Mỗi đơn vị trong bộ lọc có giá $u$. Trước khi gọi GPS supplier, cần còn1 đơn vị nếu chưa có neo và2 đơn vị nếu đã có neo. Chi phí thực trên nhánh đã bảo vệ là1 cho lần đầu,1 cho giữ neo và2 cho thử rồi tạo mới. Không đủ dự toán thì không gọi supplier; bộ chọn tiếp tục từ lịch sử đã bảo vệ.\footnote{Bộ lọc và cap khớp: \path{benchmark/engines/filtered_cover.py}, \path{benchmark/engines/matched_filter.py}. Lịch đọc: \path{benchmark/engines/paced_guard.py}. Cấp phiên và supplier: \path{core/session_budget.py}.}

\begin{table}[htbp]\centering\small
\caption{Tham số của cấu hình Geo-I / REM được giữ trong vòng L30.}
\label{tab:current-method-config}
\begin{tabularx}{\textwidth}{P{4.6cm}Y}\toprule
Đại lượng & Giá trị và ý nghĩa\\\midrule
Epoch / slot & $[0,12000)$ s; $N=8$ slot công khai; không tự reset theo GPS.\\
$C$, $H$, $U$ & $0{,}23\,\mathrm m^{-1}$; $H=12$; $U=23$ đơn vị mỗi phiên.\\
$u$ & $0{,}00125\,\mathrm m^{-1}$ cho REM và phép thử tái dùng.\\
Cap phiên / ngân sách danh nghĩa & $0{,}02875\,\mathrm m^{-1}$ / $0{,}03\,\mathrm m^{-1}$.\\
Cap epoch / danh nghĩa & $0{,}23\,\mathrm m^{-1}$ / $0{,}24\,\mathrm m^{-1}$.\\
Lịch đọc riêng tư & Ít nhất60 s giữa hai lần supplier được phép đọc.\\
Bộ chọn / chuyển động & Bảng nearest $L_{\mathrm{plan}}=10$; slack $\sigma=0{,}03$; tối đa3 lần thay một track.\\
Truy hồi / đáp án & $K=5$ Q; $L=30$ mỗi loại mỗi Q; local $k=5$.\\\bottomrule
\end{tabularx}
\end{table}

\textbf{$H$ không phải số lần đọc tối đa của bộ lọc hiện tại.} Sau lần đầu, nhánh giữ chỉ tốn1 đơn vị nhưng vẫn phải dự toán2. Với $H\geq2$, một chuỗi liên tục giữ có thể có tới $U-1=2H-2$ lần đọc: $H=12$ cho tối đa22 lần, nếu lịch có đủ sự kiện. Vòng mô phỏng600 s với khoảng60 s có không quá11 lần đọc cho cơ chế. Khi cần một ceiling đúng12 lần đọc, phải khai báo thêm ràng buộc công khai; đó chưa phải cấu hình hiện tại.

Ví dụ, tạo neo đầu tiêu $u=0{,}00125$; một lần giữ làm tổng thành $2u$; một lần thử rồi tạo mới làm tổng thành $4u=0{,}005$. Không lấy số lần tạo neo mới nhân $u$ để bỏ chi phí thử. Slot được dành trước khi engine đọc GPS, không hoàn lại phần dư từ nhánh bí mật. Phiên vượt8 slot bị từ chối, không đọc GPS hoặc gửi Q; phiên đã nhận nhưng hết cap vẫn gửi Q dự đoán. Đây là hai trạng thái khác nhau.

Sổ SQLite giữ reservation qua khởi động lại; lỗi hoặc crash không trả lại slot. Các luồng RNG được tách theo epoch, slot và chức năng bằng khóa riêng, rồi dùng PRNG của simulator. Đây là invariant vận hành dưới giả định thiết bị tin cậy, không là cơ chế chống chủ thiết bị xóa/rollback sổ. Nhiều epoch, ledger hay thiết bị của cùng người phải cộng cap hoặc dùng bộ quản lý chung. Ngân sách tọa độ không ẩn danh account/IP và không tự giải quyết S4.

\section{Ước lượng từ thông tin đã bảo vệ và chọn Q}
\label{sec:current-belief-planner}

Layer ước lượng dùng phân bố $b_t(\ell)$ trên một lưới trạng thái đường công khai để biểu diễn các vị trí còn có thể phù hợp với lịch sử neo. Nó không đọc thêm GPS để ``sửa'' ước lượng. Khi có lần đọc được phép, phân bố dự đoán được cập nhật bằng likelihood của REM/tái sử dụng; giữa các lần đọc chỉ áp dụng mô hình chuyển tiếp công khai:
\[
 \bar b_t=b_{t-1}T_{\Delta t},\qquad
 b_t(\ell)\ \propto\ \bar b_t(\ell)\,\mathcal L(Z_t\mid \ell,Z_{\mathrm{prev}}).
\]
$\mathcal L$ giữ mẫu số REM trên toàn miền, bội số của các trạng thái trùng tọa độ và khối xác suất tái sử dụng. Tuy nhiên, lưới hữu hạn, prior theo ô và mô hình chuyển tiếp là xấp xỉ; bộ lọc không điều kiện hóa đầy đủ mọi thông tin trong lịch sử nhánh. Vì vậy $b_t$ là \emph{ước lượng phục vụ utility}, không phải posterior của người dùng hoặc đối thủ đã được hiệu chuẩn. Bayesian remapping để cải thiện utility là hướng kế thừa~\cite{chatzikokolakis2017utility}.\footnote{Likelihood và chuyển tiếp: \path{benchmark/anchor_belief.py}; cờ cập nhật theo lần đọc: lớp \texttt{\_FilterBelief} trong \path{benchmark/engines/filtered_cover.py}; prior/lưới vòng hiện tại: \path{experiments/future_sumo_eval.py}.}

Mỗi track giả $j$ giữ trạng thái đường trước $q_{t-1,j}$. Miền khả thi là
\begin{equation}
 \mathcal A_{t,j}=\{v\in V_c:T_G(q_{t-1,j},v)\leq t-t_{\mathrm{prev}}\},
 \label{eq:current-feasible}
\end{equation}
trong đó $V_c$ là miền đường công khai có thể duy trì chuyển động; lần đầu dùng $V_c$. $T_G$ lấy đường đi có hướng với tốc độ free-flow công khai. Catalogue hiện tại chia trạng thái không quá40 m, tốc độ không quá8 m/s. Không nối hai làn chỉ vì tọa độ gần nhau. Ràng buộc này loại bước nhảy không có đường nối theo mô hình, nhưng chưa chứng minh phát lại đúng đèn tín hiệu, hàng xe, gia tốc hoặc đổi làn trong giao thông thật.

Từ $b_t$, bộ chọn giữ trọng số không âm $w_t(p)$ cho POI thuộc các reference nearest top5 công khai. Với $S_{10}(v)$ là hợp POI trả theo bảng planner tại $v$, mục tiêu hiện tại là
\begin{equation}
 F_t(Q)=\sum_p w_t(p)\,\mathbf1\{p\in\bigcup_{q\in Q}S_{10}(q)\},
 \qquad q_j\in\mathcal A_{t,j}.
 \label{eq:current-cover}
\end{equation}
Đây là surrogate độ phủ nearest theo belief, không phải Recall thật của bốn purpose. Mỗi trạng thái ước lượng chia trọng số đều giữa các loại có reference, rồi giữa các POI trong reference; các trạng thái không có reference không được gán đáp án giả.

Greedy xét lợi ích biên trên mọi track chưa chọn, sau đó thực hiện tối đa3 lần thay một track nếu toàn mục tiêu tăng. Các trạng thái có cùng \emph{thứ tự} signature POI được gộp bằng tie-break công khai để giảm số phép tính. Bước tiến tới mục tiêu chọn đại diện trong miền tới được; bước slack cho phép mất tối đa $\sigma=0{,}03$ điểm của $F_t$ ở \emph{cùng sự kiện và lịch sử}. Slack không là0,03 xác suất lộ GPS hay3\% trần sai số Recall. Không có ràng buộc $K$ tọa độ phải khác nhau.\footnote{Mã chọn: \path{benchmark/engines/service_cover.py}, \path{benchmark/engines/fair_cover.py}, \path{benchmark/engines/quotient_cover.py}, \path{benchmark/engines/progress_cover.py}, \path{benchmark/engines/slack_progress.py}. Factory control hiện tại: \path{benchmark/engines/public_service_planner.py}.}

Tính chất độ phủ, cận greedy và ý nghĩa của slack được chứng minh ở Mục~\ref{sec:formal-service-cover}. Cận đó áp dụng cho surrogate tại cùng một sự kiện, không phải utility của toàn quỹ đạo. Các vòng tối ưu nhiều purpose/risk không qua gate vẫn là thử nghiệm phát triển, không được nhận làm cải tiến đã xác nhận của cấu hình hiện tại.

\section{Truy hồi cố định và xếp hạng mục đích riêng tư}
\label{sec:current-local-purpose}

Mỗi Q gửi cùng danh sách \emph{mọi loại POI công khai} và $L$ cố định. Máy chủ lấy tối đa $L$ POI theo khoảng cách đường cho mỗi loại tại Q. Nó không nhận bốn truy vấn theo bốn mục đích riêng: bốn mục đích được thực hiện \emph{trên cùng ứng viên đã nhận tại thiết bị}. Sau khi hợp và bỏ trùng POI, local ranking dùng trạng thái ứng dụng tại GPS thật, category và nhu cầu thật:
\begin{table}[H]\centering\small
\caption{Bốn purpose hiện thực ở local ranking. $D_G$ là khoảng cách đường, $T_G$ là thời gian free-flow; $a$ là đích riêng.}
\begin{tabularx}{\textwidth}{P{3.3cm}Y}\toprule
Mục đích & Điều kiện / điểm xếp hạng của POI $p$\\\midrule
Gần nhất & $D_G(x,p)$.\\
Đến nhanh nhất & $T_G(x,p)$; chưa gồm congestion thời gian thực.\\
Trong bán kính $r$ & Chỉ giữ $D_G(x,p)\leq r$, rồi xếp theo $D_G$.\\
Ít vòng tới đích & $D_G(x,p)+D_G(p,a)-D_G(x,a)$; cần các đường đi hữu hạn.\\\bottomrule
\end{tabularx}
\end{table}
Tie được phá bằng ID POI công khai. Có thể trả ít hơn5 hoặc rỗng nếu không có ứng viên hợp điều kiện. Bốn purpose không bao phủ mọi nhu cầu LBS như giá, giờ mở cửa, công suất hay chỗ trống; những nhu cầu đó cần dữ liệu và giao thức bổ sung.

Trong dịch vụ tĩnh, phản hồi $L=20$ là prefix của cùng thứ tự khi $L=30$. Giữ nguyên Q, ranking và catalogue thì pool lớn chứa pool nhỏ; reference Recall từ pool lớn không thấp hơn. Đây là tính chất thu hồi ứng viên, không là privacy tăng hoặc thuật toán Q mới. Nó phải được đọc cùng chi phí phản hồi và đối chứng được cấp cùng $L$/bandwidth.

Client chỉ dùng POI đã nhận còn hiệu lực. API có cache theo epoch phản hồi; thử nghiệm L30 còn báo riêng cumulative cache của cùng phiên bản catalogue trong public epoch. Hai chính sách này không được coi là cùng TTL dịch vụ động. Khi một POI đổi trạng thái mà không có thông tin mới, trạng thái phải là chưa biết, không suy nó vẫn khả dụng. Lọc phản hồi/cache chỉ giữ lập luận hậu xử lý cho mạng khi đáp án thiếu không kích hoạt request hoặc thay lịch theo nhu cầu riêng; thiết bị có thể không trả đáp án và chờ lần refresh công khai. Lợi ích cache là hậu xử lý cục bộ có thể áp cho các phương pháp khác; current replies là phép đo chính.

\textbf{Ranh giới sensor của thực nghiệm.} Lịch60 s giới hạn supplier của Geo-I, không mọi lần ứng dụng có thể đọc vị trí cục bộ. UtilityEvaluator xếp hạng bằng GPS ground truth synthetic tại \emph{mỗi public event} và dùng đích thật làm local oracle. Việc replay L30 không đọc thêm supplier cho Q là đúng, nhưng không chứng minh utility đó đạt bằng sensor fix60 s hoặc tiêu thụ năng lượng GPS tương ứng. Một triển khai phải đo chất lượng khi vị trí local cũ, chi phí định vị và dữ liệu đích có sẵn thực tế.

\subsection{Trình tự xử lý một sự kiện công khai}
\label{sec:current-execution-order}

Các thành phần trên được gọi theo một thứ tự cố định. Trạng thái giữa hai
sự kiện gồm slot đã dành, số đơn vị đã chi, thời điểm đọc riêng tư gần nhất,
neo $Z$, phân bố $b$ và vị trí trước của từng track Q. Với một thời điểm gửi
công khai $t$, có thể đọc thuật toán theo bảy bước sau.

\begin{enumerate}
  \item \textbf{Kiểm tra phiên và thời gian.} Phiên phải được dành slot trước
  khi chạy engine; clock phải tăng và nằm trong epoch đã khai báo. Không còn
  slot thì không gọi GPS supplier và không gửi Q.
  \item \textbf{Quyết định có quan sát GPS cho cơ chế bảo vệ.} Đặt dự toán
  $a=1$ nếu chưa có neo, ngược lại $a=2$. Chỉ gọi supplier khi đã qua ít nhất
  60 s từ lần đọc riêng tư gần nhất và $\mathrm{spent}+a\le U$. Lần đầu được
  phép đọc ngay. Quyết định này không dùng GPS mới hoặc purpose thật.
  \item \textbf{Bảo vệ quan sát được phép.} Lần đầu lấy mẫu REM để tạo $Z$ và
  chi một đơn vị. Các lần sau thử $d_E(x_t,Z)+\eta\le200$ m, với
  $\eta\sim\operatorname{Lap}(1/u)$. Giữ neo thì chi một đơn vị; thử không
  đạt rồi tạo neo mới thì chi hai đơn vị. Nếu bước 2 không cho phép đọc, giữ
  neo và không chi thêm. Cập nhật thời điểm đọc chỉ khi đã đọc.
  \item \textbf{Cập nhật ước lượng.} Dự đoán $b$ theo khoảng thời gian đã qua.
  Cập nhật likelihood khi bước 3 đã đọc riêng tư, kể cả phép thử có nhiễu
  cho kết quả giữ nguyên $Z$. Event không đọc GPS chỉ dự đoán; không coi
  một neo lặp lại ở event đó là một phép đo độc lập.
  \item \textbf{Chọn và gửi Q.} Từ $b$, graph và lịch sử Q, lập miền tới được
  của từng track, chọn phủ POI, cải thiện hữu hạn rồi áp dụng slack. Gửi
  $K=5$ Q theo clock công khai, cùng mọi category và $L=30$; không gửi $Z$
  hoặc purpose. Hết cap ở một phiên đã nhận vẫn tiếp tục bước này bằng lịch
  sử đã bảo vệ.
  \item \textbf{Nhận dữ liệu dịch vụ.} Hợp phản hồi và bỏ trùng ID. Khi có
  trạng thái động, chỉ dùng bit nhận được còn hiệu lực; thông tin hết hạn
  trở thành chưa biết. Thiếu đáp án không kích hoạt request theo nhu cầu
  riêng; lần gửi sau vẫn theo clock đã định.
  \item \textbf{Trả lời tại thiết bị.} Xếp hạng ứng viên theo vị trí local
  $\hat x_t^{\mathrm{local}}$ và purpose, radius hoặc đích người dùng đã có.
  Trả tối đa năm POI phù hợp. Bước này không cập nhật $Z$, $b$, Q hoặc sổ
  ngân sách, và kết quả local không được gửi lại nhà cung cấp.
\end{enumerate}

Có hai nguồn vị trí cần phân biệt. $x_t$ ở bước 3 là quan sát đi vào cơ chế
riêng tư, còn $\hat x_t^{\mathrm{local}}$ ở bước 7 phục vụ người dùng tại
thiết bị. Vòng xác nhận tĩnh đã giả định vị trí local chính xác ở mỗi event.
Phép thử local sensor ở Chương~\ref{ch:experiments} chỉ thay bước xếp hạng
bằng vị trí giữ lại hoặc ước lượng từ fix cũ; không chấm lại Geo-I với GPS
đầu vào nhiễu. Giới hạn supplier và số fix của mô hình local không được cộng
hoặc đồng nhất thành một phép đo GNSS của điện thoại.

\textbf{Nguồn chi phí tính toán.} REM phải xét toàn miền đường cố định;
không dùng bán kính riêng quanh GPS để giảm việc này. Với $|V|$ trạng thái,
lấy mẫu neo dùng một lượt tính điểm trên $|V|$; emission belief tính trên
lưới hữu hạn, còn dự đoán dùng ma trận chuyển tiếp thưa. Bộ chọn Q xét các
miền tới được và signature POI, kèm một số lần thay track đã giới hạn.
Local ranking lấy chi phí đường có hướng và sắp các ứng viên theo từng mục
đích. Mã có cache cho context, transition và đường đi, nên chi phí lần đầu
và lần có cache khác nhau. Những mô tả này chỉ chỉ ra nơi phát sinh chi phí;
chưa là benchmark latency, CPU/RAM hay năng lượng trên thiết bị thực.


% Ideal-kernel proof for the CURRENT matched, paced REM/noisy-reuse pipeline.
% Integrator may replace the shorter privacy proof in current_method.tex.
% No production code, sampler or sealed evidence is changed by this module.
\section{Chứng minh riêng tư cho cơ chế lý tưởng}
\label{sec:formal-current-privacy}
\label{sec:current-ideal-claim}

Phần này chứng minh rằng trong cùng một epoch công khai, thay tọa độ đầu vào
bằng một chuỗi gần chúng chỉ thay xác suất quan sát máy chủ trong một hệ số
đã chặn. Đây là áp dụng của riêng tư theo khoảng cách và cơ chế dự đoán có quản lý
ngân sách~\cite{andres2013geoi,chatzikokolakis2014predictive}, không phải
một nguyên lý riêng tư mới. Chứng minh cho kernel lý tưởng, chưa là chứng chỉ pure Geo-I
cho mã float của simulator.

Ý chính là: hai đầu vào gần nhau phải tạo ra phân phối quan sát gần nhau. Bảo đảm này phải còn đúng sau kiểm tra giữ/làm mới, qua nhiều phiên và khi chuyển neo thành tập Q. Phần tiếp theo chứng minh tính khả thi trên đường, ranh giới nội dung truy vấn và điều kiện phục hồi top-$k$; Chương~\ref{ch:experiments} mới đánh giá chất lượng thực tế.

\subsection{Miền so sánh và các giả định}
\label{sec:formal-privacy-assumptions}

Giữ ngữ cảnh $\mathcal C$ cố định: graph và support công khai $V$ hữu hạn,
không rỗng; phép chiếu $\Pi$ cố định; catalogue/prior/model công khai; các
tham số $N,H,C,\theta,K,L$; epoch, thời điểm mở/đóng phiên và danh sách sự
kiện công khai hữu hạn. Giả sử $N,H$ là số nguyên dương, $H\geq1$ và
$C>0$. Với tọa độ GPS hợp lệ $x,x'$, đặt
\begin{equation}
 d_E(x,x')=\|\Pi(x)-\Pi(x')\|_2,\qquad
 D_\infty(X,X')=\max_{s,t}d_E(x_{s,t},x'_{s,t}).
 \label{eq:formal-trace-metric}
\end{equation}
$X,X'$ có cùng danh sách phiên/sự kiện, chỉ khác các tọa độ. Khoảng cách có
đơn vị m trong phép chiếu; nó không phải khoảng cách đường có hướng hoặc
định lý cho khoảng cách geodesic trên toàn cầu. Định lý bảo vệ tọa độ GPS được đưa vào cơ chế, chưa tự chuyển thành bảo đảm cho vị trí vật lý thật khi cảm biến sai. Tọa độ ở event không được
đọc vẫn thuộc $X$ để giữ cùng miền so sánh, nhưng không đi vào tính toán
riêng tư tại event đó.

\noindent\textbf{Định nghĩa (Geo-I).} Một cơ chế ngẫu nhiên $M$ thỏa $\varepsilon$-Geo-I nếu với mọi $x,x'$ và mọi biến cố đầu ra $A$,
\begin{equation}
 \Pr[M(x)\in A]\leq e^{\varepsilon d_E(x,x')}\Pr[M(x')\in A].
 \label{eq:formal-geoi-definition}
\end{equation}
Bất đẳng thức cũng áp dụng khi đổi $x,x'$. Nó chặn mức thay đổi xác suất, không bắt mọi điểm đầu ra phải nằm trong một bán kính cố định. $\varepsilon$ nhỏ hơn cho cận chặt hơn ở cùng khoảng cách~\cite{andres2013geoi}.

Chứng minh dùng các điều kiện sau:
\begin{enumerate}
 \item Tọa độ/tham số hợp lệ, hữu hạn. Support, phép chiếu và tham số không
 lấy từ GPS chưa bảo vệ; không cắt support hoặc loại mẫu theo GPS thật.
 \item REM và Laplace được lấy mẫu đúng trong số học thực với randomness
 độc lập lý tưởng. Ngẫu nhiên cho hậu xử lý độc lập với ngẫu nhiên REM/phép thử và chỉ dùng
 transcript đã bảo vệ. Đây là giả định toán học, không phải tính chất đã
 chứng nhận của NumPy PRNG/HMAC trong mã chạy.
 \item Lần đầu tạo neo bằng REM. Các lần sau dùng neo đã bảo vệ gần nhất
 làm tham chiếu, thử Laplace rồi giữ hoặc tạo neo mới. Quyết định cho phép
 đọc được tính \emph{trước} gọi GPS từ đồng hồ công khai và sổ chi phí của
 lịch sử đã bảo vệ. Bộ lọc dự toán đủ phí của nhánh đắt nhất.
 \item Sổ chung cấp tối đa $N$ slot cho chủ thể liên kết, dành trước dùng và
 không hoàn lại phần dư theo nhánh. Thiết bị tin cậy không xóa/rollback hoặc
 mở sổ thứ hai để tự cấp lại cap. Lỗi/crash quan sát được không tạo tín hiệu
 mới từ GPS/nhu cầu chưa bảo vệ.
 \item Quan sát được chứng minh là bản tin logic và phản hồi của máy chủ.
 Mọi Q, nội dung request và hành vi gửi sau cơ chế chỉ dùng $\mathcal C$
 hoặc transcript đã bảo vệ. GPS cục bộ, mục đích, category, radius, đích,
 đáp án local, click và khóa RNG không được gửi hoặc dùng để kích hoạt
 request/retry phụ thuộc bí mật.
\end{enumerate}

Thời điểm mở/đóng ứng dụng và metadata định danh được giữ cùng nhau khi so
sánh $X,X'$; định lý không bảo vệ sự thay đổi của những trường đó. Một lịch
logic cố định cũng chưa chứng minh latency của packet thực độc lập GPS.

\subsection{Bổ đề REM: phải tính cả hệ số chuẩn hóa}
\label{sec:formal-rem}

Gọi $a_v$ là tọa độ phẳng công khai của trạng thái $v\in V$, và $u>0$ là
giá của một đơn vị ngân sách, có đơn vị $\mathrm m^{-1}$. Kernel REM là
\begin{equation}
 R_u(v\mid x)=\frac{\exp[-u\|\Pi(x)-a_v\|_2/2]}{Z_u(x)},\qquad
 Z_u(x)=\sum_{w\in V}\exp[-u\|\Pi(x)-a_w\|_2/2].
 \label{eq:formal-rem-kernel}
\end{equation}
Do $V$ hữu hạn, không rỗng, mọi số hạng dương nên $0<Z_u(x)<\infty$ và mọi
trạng thái đều có xác suất dương. Với $d=d_E(x,x')$, bất đẳng thức tam giác
cho
\begin{equation}
 e^{-ud/2}\leq
 \frac{\exp[-u\|\Pi(x)-a_v\|_2/2]}{\exp[-u\|\Pi(x')-a_v\|_2/2]}
 \leq e^{ud/2}.
 \label{eq:formal-rem-score}
\end{equation}
Cộng cùng bất đẳng thức theo $v$ cho
$e^{-ud/2}\leq Z_u(x)/Z_u(x')\leq e^{ud/2}$. Vì vậy
\begin{equation}
 e^{-ud}\leq\frac{R_u(v\mid x)}{R_u(v\mid x')}\leq e^{ud}.
 \label{eq:formal-rem-ratio}
\end{equation}
Đây là $u$-Geo-I theo $d_E$. Hệ số $1/2$ trong số mũ cần thiết cho cách
chặn này: một nửa trả cho score, một nửa cho chuẩn hóa. Cộng xác suất của
các trạng thái cùng tọa độ vẫn giữ cận; không cần giả định mỗi vertex có
tọa độ khác nhau. Support cố định là điều kiện riêng tư, còn khả thi trên
đường của Q thuộc bước hậu xử lý sau REM.\footnote{Kernel hiện tại:
\path{core/mechanisms.py}, lớp \texttt{RoadExponential}.}

\subsection{Bổ đề phép thử: cả giữ và không giữ đều phải riêng tư}
\label{sec:formal-private-test}

Tại một lịch sử đã bảo vệ cố định, tham chiếu $h$ là cùng một điểm dưới
$X,X'$. Đặt $g_h(x)=\|\Pi(x)-h\|_2$. Hàm này là 1-Lipschitz:
$|g_h(x)-g_h(x')|\leq d_E(x,x')$. Với
$Y=g_h(x)+\eta$, $\eta\sim\operatorname{Lap}(1/u)$, mật độ là
\begin{equation}
 f_x(y)=\frac{u}{2}e^{-u|y-g_h(x)|},\qquad
 \frac{f_x(y)}{f_{x'}(y)}
 \leq e^{u|g_h(x)-g_h(x')|}\leq e^{ud_E(x,x')}.
 \label{eq:formal-laplace-density}
\end{equation}
Tích phân trên hai miền $(-\infty,\theta]$ và $(\theta,\infty)$, với
$\theta$ hữu hạn công khai, cho
\begin{equation}
 \begin{split}
 q_x&=\Pr[Y\leq\theta],\\
 q_x&\leq e^{ud_E(x,x')}q_{x'},\qquad
 1-q_x\leq e^{ud_E(x,x')}(1-q_{x'}).
 \end{split}
 \label{eq:formal-test-two-branches}
\end{equation}
Đổi $x,x'$ cho hai cận ngược. Cả $q_x$ và $1-q_x$ đều dương với các
tọa độ hữu hạn. Phép thử không phải rule ``dưới ngưỡng luôn giữ, trên ngưỡng
luôn tạo mới''; cả hai kết quả đều được trả phí riêng tư.\footnote{Phép thử:
\path{core/mechanisms.py}, lớp \texttt{PrivateReuseSMREM}.}

Với khoảng cách $d=\|\Pi(x)-h\|_2$, xác suất giữ có biểu thức cụ thể
\begin{equation}
 q(d)=\begin{cases}
 1-\tfrac12 e^{-u(\theta-d)},&d\leq\theta,\\
 \tfrac12 e^{-u(d-\theta)},&d>\theta.
 \end{cases}
 \label{eq:formal-hold-probability}
\end{equation}
Tại $u=0{,}00125\,\mathrm m^{-1}$, $\theta=200$ m, khoảng cách 50 m cho
$q\simeq0{,}5855$, còn 1.200 m cho $q\simeq0{,}1433$. Ngưỡng 200 m
không là cận cứng của sai số. Với $r\geq\theta$, điều kiện theo lịch sử trước
test rồi lấy trung bình cho
\begin{equation}
 \Pr[\text{giữ}\ \text{và}\ d>r]
 \leq\tfrac12e^{-u(r-\theta)}.
 \label{eq:formal-far-hold-joint}
\end{equation}
Đây là cận biến cố chung tại một event, không phải
$\Pr[d>r\mid\text{giữ}]$; xác suất có điều kiện đó có thể bằng 1 khi
tham chiếu luôn xa. Không chuyển cận này thành bảo đảm sai số của Q hoặc
Recall người dùng.

\subsection{Chi phí theo nhánh là cận của xác suất chung}
\label{sec:formal-joint-branch}

Để chứng minh, xét transcript mở rộng $E$: tại mỗi event nó ghi bỏ qua,
tạo neo đầu, giữ neo hoặc tạo neo mới, kèm ID trạng thái khi lấy mẫu REM.
$E$ không chứa GPS thật, giá trị nhiễu $\eta$, các giá trị ngẫu nhiên gốc hoặc khóa RNG;
nó là công cụ chứng minh, không yêu cầu công bố các trường này. Gọi
$e_{<i}$ là phần lịch sử trước sự kiện $i$ (prefix) cố định; neo $h$, số đơn vị đã chi và đồng hồ đọc
đều được xác định từ prefix đó. Tại event được đọc sau lần đầu, kernel
\emph{chung} của nhánh và neo là
\begin{equation}
 \begin{split}
 K_x(\text{giữ}\mid e_{<i})&=q_x,\\
 K_x(\text{mới},v\mid e_{<i})&=(1-q_x)R_u(v\mid x).
 \end{split}
 \label{eq:formal-extended-step}
\end{equation}
Fresh REM dùng randomness độc lập với phép thử, nên có tích thứ hai.
Bổ đề trên cho cận một đơn vị $u$ ở nhánh giữ và hai đơn vị $2u$ ở nhánh
mới. Lần đầu không có test, chỉ tốn $u$. Event bị bỏ qua theo pacing/cap
không đọc GPS, kernel không phụ thuộc $x$ và tốn 0.

\textbf{Không được bỏ xác suất nhánh hoặc chuẩn hóa lại đường đã chọn.}
Các cận này so sánh xác suất chung của \emph{cùng} outcome khi điều kiện
theo \emph{cùng prefix}, không phải định lý rằng một cơ chế được chọn sau
chạy có epsilon bằng số lần tạo neo mới. Nếu fresh REM lại lấy đúng $h$,
outcome ``mới, $h$'' vẫn khác outcome ``giữ'' trong $E$. Khi chỉ quan sát
neo, phải cộng cả hai khối xác suất; không suy nhánh từ việc hai tọa độ bằng
nhau. Không cần công bố, hoặc giả định đối thủ biết, nhánh nội bộ để lấy
biên này.

\subsection{Định lý bộ lọc và hợp thành qua phiên}
\label{sec:formal-adaptive-cap}

Đặt $U=2H-1$, $u=C/(NU)$ như bộ lọc khớp cap hiện tại. Trước đọc đầu cần
còn ít nhất 1 đơn vị; các lần sau cần ít nhất 2. Nếu không đủ dự toán, không
đọc GPS và tiếp tục hậu xử lý. Với mọi đường $e$, gọi
$c_i(e)\in\{0,1,2\}$ là chi phí tại event $i$. Quy tắc dự toán cho, bằng
quy nạp,
\begin{equation}
 \sum_{i\text{ trong phiên }s}c_i(e)\leq U,
 \qquad \sum_i c_i(e)\leq NU.
 \label{eq:formal-pathwise-cap}
\end{equation}
Chi phí có thể thích nghi theo nhánh; cận trên là cho \emph{mọi đường},
không phải chi phí kỳ vọng. Bộ lọc dùng integer units để hiện thực invariant
này, nhưng chỉ invariant đó chưa chứng minh chất lượng sampler.

\noindent\textbf{Định lý (epoch lý tưởng).} Dưới các giả định đã nêu, với
mọi $X,X'$ cùng $\mathcal C$ và mọi biến cố quan sát máy chủ $A$,
\begin{equation}
 \Pr[\mathsf O(X)\in A]
 \leq e^{C D_\infty(X,X')}\Pr[\mathsf O(X')\in A].
 \label{eq:current-transcript-bound}
\end{equation}

\noindent\emph{Chứng minh.} Đánh số hữu hạn tất cả event trong epoch theo
thứ tự công khai. Tại cùng $e_{<i}$, bộ lọc có cùng số dư, tham chiếu và
thời điểm đọc trước dưới $X,X'$; admission, dự toán và pacing do đó cho
cùng quyết định. Không điều kiện theo hai prefix khác nhau để so kernel.
Ở event có lấy mẫu, REM và hai nhánh Laplace có support dương chung. Ở
event không lấy mẫu, outcome xác định giống nhau. Vì vậy các đường khả thi
có cùng support dưới hai input; không xuất hiện một đường chỉ có xác suất
dương ở một phía do bộ lọc.

Cùng prefix không có nghĩa hai lần chạy thực tế luôn đi cùng nhánh: ta
đang chặn xác suất của từng đường tương ứng, rồi cộng các đường.

\noindent
Áp chain rule cho xác suất \emph{chung} của một đường $e$ và các cận từng
event:
\begin{equation}
 \begin{split}
 p_X(e)&=\prod_i K_{x_i}(e_i\mid e_{<i}),\\
 \frac{p_X(e)}{p_{X'}(e)}
 &\leq \exp\!\left(u\sum_i c_i(e)d_E(x_i,x'_i)\right)\\
 &\leq \exp\!\left(uNU D_\infty(X,X')\right)
 =e^{C D_\infty(X,X')}.
 \end{split}
 \label{eq:formal-path-ratio}
\end{equation}
Các quyết định dừng theo cap đã nằm trong chain rule qua prefix, nên
không cộng thêm một ``phí dừng'' hoặc giả định stopping time độc lập input.
Điều cần là nó chỉ dùng lịch công khai và $E$, cùng với dự toán trước đọc.
Cộng bất đẳng thức cho mọi đường trong một biến cố cho cận của $E$. Bước
hậu xử lý bên dưới chuyển cận sang $\mathsf O$. \hfill$\square$

Đây là áp dụng của bounded predictive composition~\cite{chatzikokolakis2014predictive};
Theorem 1 của nguồn gốc dùng mức chi lớn nhất trên các run.\footnote{Bản tác
giả: \url{https://arxiv.org/abs/1311.4008}; Theorem 1 và chứng minh trong phụ
lục. Bộ lọc hiện tại: \path{benchmark/engines/filtered_cover.py},
\path{benchmark/engines/matched_filter.py},
\path{benchmark/engines/paced_guard.py}; sổ qua phiên:
\path{core/session_budget.py}.}

\noindent\textbf{Hệ quả (chỉ thay một tọa độ đầu vào).} Nếu $X,X'$ giống nhau tại mọi sự kiện trừ $i_0$, đặt $r=d_E(x_{i_0},x'_{i_0})$. Trong (\ref{eq:formal-path-ratio}), mọi khoảng cách khác bằng 0 và $c_{i_0}(e)\leq2$, nên toàn quan sát máy chủ thỏa
\begin{equation}
 \Pr[\mathsf O(X)\in A]\leq e^{2ur}\Pr[\mathsf O(X')\in A].
 \label{eq:formal-single-coordinate}
\end{equation}
Lấy biên và hậu xử lý giữ cận như trên, kể cả các sự kiện sau thay đổi phản hồi do lịch sử neo. Với $r=100$ m và $u=0{,}00125\,\mathrm m^{-1}$, hệ số là $e^{0{,}25}\simeq1{,}284$. Đây là so sánh hai trace chỉ khác một mẫu GPS; hai tuyến khác nhau ở nhiều thời điểm không được dùng cận này. Nó không tự bảo vệ thuộc tính nhà/nơi làm việc được suy từ cả tuyến hoặc giờ bật ứng dụng.

Mỗi slot có cap hiệu lực $Uu=C/N$; engine nhận ngân sách danh nghĩa $2Hu$,
không phải cap này. Tham số hiện tại $N=8,H=12,C=0{,}23\,\mathrm m^{-1}$
cho $U=23$, $u=0{,}00125\,\mathrm m^{-1}$ và cap phiên
$0{,}02875\,\mathrm m^{-1}$. $H$ không phải ceiling số lần GPS: đường giữ
liên tục có thể đọc $2H-2$ lần khi $H\geq2$ và lịch đủ dài; $H=1$ có một
lần đầu. Supplier pacing 60 s trong phiên 600 s giới hạn ở 11 lần bằng
đồng hồ riêng. Ví dụ, ba lần đọc lần lượt tạo neo đầu, giữ neo, rồi tạo mới dùng $(1+1+2)u=0{,}005\,\mathrm m^{-1}$. Việc sinh $K=5$ Q và nhận $L=30$ POI mỗi loại từ lịch sử đó không tự cộng thêm phí Geo-I: chúng là hậu xử lý theo các giả định đã nêu. Thay $K,L$ vẫn có thể đổi độ chính xác, chi phí hoặc khả năng suy luận thực nghiệm. Các epoch/sổ khác của cùng chủ thể hợp thành tiếp; reset
không cho một cap chung mới cho toàn lịch sử.

\subsection{Hậu xử lý: từ neo nội bộ tới toàn quan sát máy chủ}
\label{sec:formal-server-postprocessing}

Điều kiện theo $E=e$, belief, các track đường, Q, thứ tự Q và request chỉ
phụ thuộc $e,\mathcal C$ và randomness hậu xử lý. Máy chủ trả lời từ Q và
ngữ cảnh/chiến lược của nó; phản hồi không phải kênh GPS thật thứ hai. Vì
vậy tồn tại một kernel chung $T_{\mathcal C}(A\mid e)$ không phụ thuộc
$X$ ngoài $e$, thỏa
\begin{equation}
 \begin{split}
 \Pr[\mathsf O(X)\in A]
 &=\sum_e p_X(e)T_{\mathcal C}(A\mid e)\\
 &\leq e^{CD_\infty(X,X')}
       \sum_e p_{X'}(e)T_{\mathcal C}(A\mid e).
 \end{split}
 \label{eq:formal-postprocessing}
\end{equation}
Kernel này bao gồm các request logic, thứ tự/độ dài đã mô hình hóa, phản
hồi và việc tiếp tục dự đoán sau khi hết cap. Máy chủ có thể thích nghi theo
những bản tin nó đã thấy và randomness riêng; chứng minh vẫn dùng cùng
chiến lược, miễn không có feedback mới từ GPS/nhu cầu chưa bảo vệ. Một
belief xấp xỉ hoặc planner chưa tối ưu vẫn giữ hậu xử lý; độ chuẩn của
posterior và utility là câu hỏi khác.

Local ranking được phép đọc GPS và nhu cầu thật vì đáp án local không nằm
trong $\mathsf O$. Nếu công bố đáp án đó, nó không còn là hậu xử lý của
$E$; không tự áp (\ref{eq:formal-postprocessing}). Lọc POI, cache và trạng
thái unknown chỉ giữ ranh giới mạng khi đáp án thiếu không sinh retry hoặc
đổi thời điểm request theo bí mật. Account/IP, packet latency thực, crash
và truy cập bộ nhớ ngoài mô hình phải được phân tích riêng.

Tính chất không thêm kênh purpose cho S7 được chứng minh riêng. Cận tọa độ
ở đây không suy intent độc lập với tuyến hoặc giờ bật ứng dụng.

\subsection{Hệ quả Bayesian: thông tin thêm do quan sát}
\label{sec:formal-bayesian-odds}

Xét hai giả thuyết $H_0:X$ và $H_1:X'$ với cùng $\mathcal C$ và prior dương,
có thể đã điều kiện theo tri thức phụ cố định của đối thủ. Với biến cố
quan sát $A$ có xác suất dương, định lý và chiều ngược cho
\begin{equation}
 e^{-CD_\infty(X,X')}
 \leq
 \frac{\Pr[H_0\mid A]/\Pr[H_1\mid A]}
      {\Pr[H_0]/\Pr[H_1]}
 =\frac{\Pr[A\mid H_0]}{\Pr[A\mid H_1]}
 \leq e^{CD_\infty(X,X')}.
 \label{eq:formal-posterior-odds}
\end{equation}
Ở đây odds là tỷ số xác suất của hai giả thuyết. Quan sát mới chỉ được nhân tỷ số trước quan sát trong khoảng này. Kết quả
không đòi prior bằng nhau, nhưng không sửa prior đã gần chắc chắn, không
chặn Hit100 ở một tỷ lệ cố định và không cho MAE tối thiểu. Với hai giả
thuyết hỗn hợp về tập trace, có thể lấy tích phân để có cùng cận $e^{Cr}$
\emph{nếu mọi cặp trace giữa hai support đều có $D_\infty\leq r$} và cùng
ngữ cảnh. Không áp cận cho hai nhóm tùy ý có tuyến rất xa nhau.

Với $r=100$ m, một REM hoặc nhánh giữ có cận $e^{ur}=e^{0{,}125}$ khi
$u=0{,}00125\,\mathrm m^{-1}$; nhánh test rồi tạo mới có cận $e^{0{,}25}$.
Nếu chỉ có hai giả thuyết với prior $1/2$--$1/2$ tại một lịch sử cố định, chỉ xét riêng một đầu ra REM (không kèm kết quả phép thử), cận posterior tối đa là $e^{0{,}125}/(1+e^{0{,}125})\simeq53{,}12\%$. Đây là ví dụ cho một quan sát với hai ứng viên, không phải độ chính xác của attacker benchmark. Cả epoch có cận $e^{Cr}=e^{23}$, rất lỏng ở thang đó. Theorem xác định cách
hợp thành, không biến epsilon thành Hit/MAE hoặc chứng minh mọi S1--S10,
identity, giờ đầu/cuối, đồng hành hay điểm nhà đã được bảo vệ. Benchmark
vẫn cần để đánh giá utility và sức tấn công thực dụng.

\subsection{Phản ví dụ khi bỏ điều kiện và giới hạn thực thi}
\label{sec:formal-counterexamples}

\textbf{Giữ neo không miễn phí.} Lấy tham chiếu công khai $h=0$, ngưỡng
$\theta=0$ và hai vị trí cách $h$ lần lượt 0 và $r>0$. Khi chỉ xét outcome
giữ của Laplace test, $q_0=1/2$ và $q_r=e^{-ur}/2$, tỷ số bằng $e^{ur}$.
Đếm riêng số neo fresh rồi gán nhánh giữ chi phí 0 bỏ qua đúng lượng tin
đã lộ. Ngưỡng 0 chỉ dùng làm phản ví dụ, không đổi ngưỡng 200 m hiện tại.

\textbf{Support theo GPS tạo xác suất 0.} Trên support công khai
$V=\{0,1\}$ của một đường thẳng, chỉ nhận vertex cách GPS không quá 0,6
sẽ cho vertex 0 xác suất dương ở $x=0{,}6-\delta$ và xác suất 0 ở
$x'=0{,}6+\delta$, với $0<\delta<0{,}1$. Hai input cách $2\delta$ vẫn có
tỷ số vô hạn. Tương tự, một quy tắc giữ neo chắc chắn dưới ngưỡng và fresh
chắc chắn trên ngưỡng có thể lộ việc qua ranh giới; nhiễu test là điều kiện
của proof, không chỉ thủ thuật utility.

\textbf{Cap không chặn kênh bên ngoài cơ chế.} Nếu gửi hoặc không gửi
request được quyết định trực tiếp từ GPS thật, hai input ở hai phía một
ngưỡng có thể tạo biến cố ``có request'' với xác suất 1 và 0, dù tổng chi
REM không vượt cap. Điều tương tự xảy ra khi retry/click phụ thuộc mục
đích riêng hoặc memory dump đưa raw GPS/khóa RNG vào quan sát. 
Mã hiện tại dùng float64, Gumbel-max/Laplace và NumPy PRNG. Ngẫu nhiên hữu hạn có thể tạo xác suất 0 phụ thuộc input, trái với support dương lý tưởng. Monte Carlo, log-domain và seed riêng không chứng minh pure Geo-I. HMAC tách luồng và sổ SQLite xác nhận invariant dưới giả định thiết bị tin cậy, chưa chứng nhận ngẫu nhiên lý tưởng hoặc chống rollback. Cần sampler/entropy có chứng minh, hoặc cận approximate được chứng nhận; hiện chưa có chứng chỉ đó.\footnote{Giới hạn sampler: \path{core/mechanisms.py}; phản ví dụ: \path{tests/test_sampler_support.py}. Nguồn: I. Mironov, CCS 2012, \url{https://www.microsoft.com/en-us/research/publication/on-significance-of-the-least-significant-bits-for-differential-privacy/}.}

% Service properties of the CURRENT legacy_l10 Q planner and local ranking.
% Additive proof module; no executable or sealed empirical evidence is changed.
% Intended for the method chapter, before empirical evaluation.
\section{Tính chất của bước chọn Q và phục hồi dịch vụ}
\label{sec:formal-current-service}

Phần này tách ba câu hỏi: Q có thể di chuyển theo mô hình đường hay không;
nhu cầu riêng có làm đổi bản tin mạng hay không; và tập POI nhận về có đủ
để trả đúng top-$k$ hay không. Các tính chất này không thay thế định lý
Geo-I hoặc chứng minh phương pháp vượt đối chứng. Chúng áp dụng cho cấu
hình hiện tại: $K=5$, bảng chọn Q có $L_{\mathrm{plan}}=10$, phản hồi
có $L=30$ và đáp án local có $k=5$.

\subsection{Khả thi theo đường có hướng}
\label{sec:formal-service-reachability}

Giữ graph công khai $G=(V,E)$ cố định. Mỗi cạnh có độ dài không âm và tốc
độ dương; $T_G(v,w)$ là thời gian đường đi ngắn nhất có hướng, bằng
$+\infty$ khi không tới được. Cạnh nối hai trạng thái trùng tọa độ có thể
có chi phí 0. Miền $V_c\subseteq V$ là thành phần liên thông mạnh công
khai không rỗng được chọn trước chuyến. Gọi $v_{t,j}$ là \emph{trạng thái
làn--độ tiến nội bộ} của track $j$, còn $g(v_{t,j})$ là tọa độ Q gửi đi.
Với lịch công khai tăng nghiêm ngặt, miền chọn là
\begin{equation}
 \mathcal A_{t,j}=\begin{cases}
 V_c,&\text{sự kiện đầu},\\
 \{v\in V_c:T_G(v_{t-1,j},v)\leq\Delta t\},&\text{các sự kiện sau},
 \end{cases}
 \qquad \Delta t=t-t_{\mathrm{prev}}>0.
 \label{eq:formal-service-buckets}
\end{equation}

\noindent\textbf{Tính chất 1 (khả thi từng track).}
Nếu mọi bước chọn và thay trạng thái giữ $v_{t,j}\in\mathcal A_{t,j}$,
mỗi track có một đường đi trong $G$ giữa hai sự kiện liên tiếp với thời
gian không quá $\Delta t$.

\noindent\emph{Chứng minh.}
Miền đầu không rỗng. Ở bước sau, trạng thái trước vẫn thuộc $V_c$ và
$T_G(v_{t-1,j},v_{t-1,j})=0$, nên miền chọn không rỗng. Greedy chọn trong
miền; gộp trạng thái chỉ lấy tập con của miền; thay một track vẫn dùng
miền của track đó. Bước tiến tới mục tiêu chọn một trạng thái tới được
có cùng signature POI; bước slack cũng chỉ xét trạng thái tới được.
Quy nạp cho tính chất. Nối các đường giữa sự kiện và cho phép chờ bù thời
gian tạo một lộ trình theo mô hình graph. \hfill$\square$

Đây là tính chất của \textbf{trạng thái nội bộ}, không phải chứng minh
mọi phép suy track từ tọa độ đều khả thi. Hai làn khác nhau có thể trùng
tọa độ; phép truy hồi POI bằng trạng thái gần tọa độ nhất có thể chọn một
làn khác. Vì vậy chỉ khẳng định tồn tại cách gán trạng thái khả thi cho
chuỗi tọa độ đã sinh. Xáo thứ tự Q không xóa khả năng đối thủ nối track
bằng hình học. Các track cũng có thể trùng trạng thái; không có bảo đảm
tránh va chạm hoặc $K$ vị trí khác nhau.

Graph dùng tốc độ free-flow và danh mục trạng thái rời rạc công khai;
vòng hiện tại dùng khoảng chia không quá40 m và tốc độ không quá8 m/s.
Điều này không xác nhận đèn tín hiệu, tắc đường, gia tốc hoặc đổi làn
trong giao thông thật. Phép tính đường ngắn nhất lý tưởng cho tính chất
trên; kiểm tra SciPy/float chỉ xác nhận mô hình số đang chạy, chưa chứng
nhận sai số đối với đường liên tục.\footnote{Miền làn công khai:
\path{data/lane_states.py}; graph có hướng và phép tới được:
\path{evaluation/lane_travel.py}; các bước giữ miền:
\path{benchmark/engines/quotient_cover.py},
\path{benchmark/engines/progress_cover.py},
\path{benchmark/engines/slack_progress.py}.}

\subsection{Mục tiêu độ phủ và ý nghĩa chính xác của slack}
\label{sec:formal-service-cover}

Tại một sự kiện và lịch sử đã bảo vệ cố định, gọi $S_{10}(v)$ là hợp các
ID POI trong signature nearest của bộ chọn tại trạng thái $v$. Bảng này
theo phép ánh xạ tọa độ của dịch vụ, không được lấy trực tiếp từ GPS thật.
Trọng số $w_t(p)\geq0$ được tính từ belief trên lưới công khai. Trong cấu
hình hiện tại, mô hình gốc tạo trọng số từ \textbf{reference nearest
top5}; wrapper giữ các trọng số đó và dùng signature top10 để tính độ
phủ. Không đồng nhất hai độ sâu này với phản hồi $L=30$.

Với một lựa chọn trạng thái cho mỗi track, đặt
\begin{equation}
 F_t(Q)=\sum_p w_t(p)\,
 \mathbf1\{p\in\textstyle\bigcup_{j=1}^K S_{10}(v_{t,j})\}.
 \label{eq:formal-service-cover}
\end{equation}
Đây là điểm độ phủ nearest theo belief xấp xỉ. Nó không phải Recall thật
của bốn purpose, cũng không là kỳ vọng đã hiệu chuẩn cho một người dùng.
Trạng thái lưới không có reference đóng góp khối trọng số 0; không tạo
POI hoặc điểm số giả để thay reference rỗng.

\noindent\textbf{Tính chất 2 (độ phủ có lợi ích biên giảm dần).}
Xem mỗi phần tử chọn là cặp có nhãn $(j,v)$, với
$v\in\mathcal A_{t,j}$. Khi đó $F_t(\varnothing)=0$, $F_t$ không giảm
và là hàm submodular. Cụ thể, nếu $A\subseteq B$ thì
\[
 \Delta(e\mid A)=\sum_{p\in S_{10}(e)\setminus S_{10}(A)}w_t(p)
 \ \geq\
 \Delta(e\mid B)\geq0.
\]
Mở rộng tập chỉ thêm POI; một POI đã được phủ không tạo lợi ích lần hai.
Ràng buộc lấy không quá một phần tử mỗi nhãn tạo một matroid phân hoạch
(partition matroid). Vì mỗi miền không rỗng và $F_t$ không giảm, có thể hoàn
thành lựa chọn thành đúng $K$ phần tử, một phần tử cho mỗi track. Hai nhãn
vẫn có thể chọn cùng trạng thái đường.

\noindent\textbf{Hệ quả (cận greedy kế thừa).}
Với oracle lợi ích biên chính xác trên toàn bộ các miền không rỗng,
greedy toàn cục đang dùng đạt $F_t(Q_{\mathrm g})\geq F_t(Q^*)/2$, trong
đó $Q^*$ tối ưu \emph{cho cùng sự kiện, belief và miền chọn}.
Đây là cận greedy chuẩn dưới ràng buộc matroid~\cite{calinescu2011matroid}, không là đóng góp
nguyên lý mới. Có thể chứng minh trực tiếp như sau. Gọi $a_i$ là phần tử
greedy chọn ở bước $i$, $S_i=\{a_1,\ldots,a_i\}$ và $o_i$ là phần tử
của $Q^*$ thuộc cùng track với $a_i$. Khi chọn $a_i$, track đó chưa được
chọn, nên $o_i$ còn khả thi. Greedy cho
$\Delta(o_i\mid S_{i-1})\leq\Delta(a_i\mid S_{i-1})$. Do tính không
giảm và submodular,
\begin{equation}
 \begin{split}
 F_t(Q^*)&\leq F_t(S_K\cup Q^*)\\
 &\leq F_t(S_K)+\sum_{i=1}^K\Delta(o_i\mid S_K)\\
 &\leq F_t(S_K)+\sum_{i=1}^K\Delta(o_i\mid S_{i-1})
 \leq2F_t(S_K).
 \end{split}
 \label{eq:formal-service-greedy-half}
\end{equation}

Gộp những \emph{signature có cùng thứ tự} giữ độ phủ và đại diện có
tie-break công khai tốt nhất tại bước hiện tại. Thay một track được nhận
chỉ khi toàn mục tiêu tăng; tiến tới mục tiêu giữ signature nên giữ
$F_t$. Giới hạn tối đa3 lần thay không chứng minh đã tới tối ưu cục bộ.
Gọi $Q_{\mathrm{pre}}$ là lựa chọn trước slack và
$\delta_F=\texttt{utility\_slack}$. Mỗi thay đổi slack được kiểm tra lại
với cùng floor $F_t(Q_{\mathrm{pre}})-\delta_F$, do đó trong số học lý
tưởng
\begin{equation}
 F_t(Q_{\mathrm{final}})\geq F_t(Q_{\mathrm{pre}})-\delta_F
 \geq\tfrac12F_t(Q^*)-\delta_F.
 \label{eq:formal-service-slack-floor}
\end{equation}
Ở bước đầu hoặc khi slack bằng0, bước slack không đổi lựa chọn. Giá trị
hiện tại $\delta_F=0{,}03$ là mất tối đa0,03 \emph{điểm surrogate cùng
bước}, không phải3\% sai số Recall hay xác suất lộ vị trí. Mã chấp nhận
floor với tolerance $10^{-12}$; đó là kiểm tra giá trị float ghi nhận,
không là chứng chỉ sai số $10^{-12}$ so với số thực. Cận $1/2$ cần oracle
chính xác, không tự áp như chứng chỉ executable. Trạng thái sau slack
làm đổi miền bước sau, nên không suy cận độ phủ cho toàn chuyến hoặc
chất lượng người dùng.\footnote{Trọng số reference5 được giữ trong
\path{benchmark/response_aware_belief.py}; nguồn resource:
\path{experiments/future_sumo_eval.py}, hàm \texttt{native\_resources}.
Greedy, độ phủ và full-action guard:
\path{benchmark/engines/service_cover.py},
\path{benchmark/engines/fair_cover.py}; factory active:
\path{benchmark/engines/public_service_planner.py}.}

\subsection{Không thêm kênh request theo nhu cầu riêng}
\label{sec:formal-service-noninterference}

Gọi $\psi$ là chuỗi nhu cầu local: purpose, category, bán kính, đích và
top-$k$. Giữ cùng trace vị trí $X$, ngữ cảnh công khai $\mathcal C$, lịch
mở/đóng, admission, sự kiện, mọi loại được request và $L$ cố định. Giữ
cùng chiến lược máy chủ. Giả sử $\psi$ chỉ đi vào local ranking, không
đổi Q, lịch đọc/gửi, expiry, reset hoặc retry. Với
$\mathsf W(X,\psi;\mathcal C)$ là transcript \emph{mạng logic}, ta có
\begin{equation}
 \mathsf W(X,\psi;\mathcal C)
 \overset{\mathrm d}{=}
 \mathsf W(X,\psi';\mathcal C)
 \quad\text{cho mọi }\psi,\psi'.
 \label{eq:current-purpose-noninterference}
\end{equation}

\noindent\emph{Chứng minh.}
Ghép cùng randomness của REM, planner, thứ tự Q và máy chủ. Ở sự kiện
đầu, các đầu vào tạo bản tin giống nhau. Nếu các prefix mạng giống nhau,
lịch sử dùng cho sự kiện tiếp theo giống nhau; nhu cầu local không sửa
chúng. Mỗi Q request mọi category công khai với cùng $L$, nên cả request
và response tiếp theo giống nhau. Quy nạp cho bình đẳng từng bản tin
trên cùng randomness, mạnh hơn bình đẳng phân bố. \hfill$\square$

API mạng không nhận \texttt{QuerySpec}; hàm trả lời local không gọi GPS
supplier của Geo-I, không fetch và không tiến đồng hồ cache. Điều kiện
này bảo vệ \emph{nội dung request có điều kiện theo cùng vị trí/lịch}.
Nó không nói intent độc lập với tuyến tới bệnh viện, giờ bật ứng dụng,
account/IP, click hoặc đặt chỗ. Logical-clock proof cũng không chứng
minh latency thực độc lập nhu cầu nếu CPU/cache contention hay lỗi local
ảnh hưởng packet. Nếu thiếu đáp án làm ứng dụng gửi request mới theo
nhu cầu thì điều kiện đã bị phá.\footnote{Ranh giới API:
\path{benchmark/query_purpose.py},
\path{benchmark/budgeted_geoi_lbs.py}; kiểm tra local scope chung cho
hai cách dùng cache: \path{benchmark/versioned_static_geoi_lbs.py}.}

\subsection{Khi nào local top-$k$ là đáp án chính xác?}
\label{sec:formal-service-local-exactness}

Giữ vị trí local đúng trong mô hình: $v$ là trạng thái gần GPS hiện tại
theo phép ánh xạ công khai. Giữ một catalogue/version và nhu cầu $\psi$.
Với category đã chọn, các điểm số hiện thực là $D_G(v,p)$ cho nearest,
$T_G(v,p)$ cho fastest, $D_G(v,p)$ với điều kiện $D_G(v,p)\leq r$ cho
within-radius, và
$D_G(v,p)+D_G(p,a)-D_G(v,a)$ cho detour. Chỉ nhận điểm số hữu hạn.
Nếu đích $a$ không tới được thì detour có reference rỗng. Bán kính là
khoảng cách đường có hướng, còn fastest dùng tốc độ free-flow, không
phải thời gian giao thông đo thực. Tie được phá bằng ID POI công khai.

Gọi $E_\psi(v)$ là miền POI hợp điều kiện và $\prec$ là thứ tự toàn phần
theo điểm số rồi ID. Đặt
\begin{equation}
 R=\operatorname{Top}_k(E_\psi(v)),\quad
 r_R=|R|=\min(k,|E_\psi(v)|),\quad
 P(A)=\operatorname{Top}_k(A\cap E_\psi(v)),
 \label{eq:formal-service-reference}
\end{equation}
trong đó $A$ là tập ID nhận hợp lệ. Với dịch vụ tĩnh hiện tại, tính hợp lệ
yêu cầu đúng catalogue/version. Nếu xét availability động, phải bổ sung
trạng thái \emph{hiện tại đã biết}; chỉ có metadata không đủ xác nhận POI
đang hoạt động. Detour cần đích local đã biết; dùng endpoint thật trong
evaluator không chứng minh thiết bị suy được đích từ quá khứ.

\noindent\textbf{Tính chất 3 (thu hồi reference).}
Khi $r_R>0$, nếu đáp án và reference dùng \emph{cùng thứ tự và miền đúng},
\begin{equation}
 R\cap P(A)=R\cap A,\qquad
 \operatorname{Recall}@k(A)=\frac{|R\cap A|}{r_R}.
 \label{eq:formal-service-recall-identity}
\end{equation}
Một $p\in R\cap A$ có hạng toàn miền không quá $k$; bỏ POI ngoài $A$
không thể làm hạng đó lớn hơn. Do đó $p$ vẫn nằm trong $P(A)$. Chiều
ngược hiển nhiên vì $P(A)\subseteq A$. Đặc biệt, $R\subseteq A$ là điều
kiện cần và đủ cho $P(A)=R$ và Recall bằng1. Không cần tải toàn catalogue,
nhưng không có bảo đảm bộ chọn Q luôn thu hồi đủ $R$.

Reference có thể ít hơn5 POI. Khi $r_R=0$, Recall là N/A và số reference
rỗng phải được báo riêng. Khi reference có nghĩa nhưng câu trả lời rỗng,
Recall bằng0; không đổi mẫu số thành kích thước reference của vị trí
ước lượng. Completion đủ số lượng cũng không đồng nghĩa Recall đúng:
có thể trả5 POI hợp điều kiện nhưng cả5 đều ngoài reference top5.
\footnote{Điểm số và tie-break:
\path{benchmark/query_purpose.py}, lớp \texttt{MultiPurposeRoadRanking};
reference/mẫu số của vòng hiện tại:
\path{experiments/qplanner_study_20261006_v2.py},
lớp \texttt{UtilityEvaluator}.}

\subsection{Tăng L: điều kiện đơn điệu và phản ví dụ}
\label{sec:formal-service-depth-monotonicity}

Giữ \emph{đúng cùng Q}, graph, phép truy cập tọa độ, catalogue tĩnh và
thứ tự dịch vụ. Với mỗi Q/category, top20 là prefix của top30, nên hợp
ID sau khử trùng thỏa $A_{20}\subseteq A_{30}$; có thể bằng nhau khi ít
POI. Nếu giữ cùng reference và local ranking đúng như Tính chất 3,
\begin{equation}
 \operatorname{Recall}@k(A_{30})-
 \operatorname{Recall}@k(A_{20})
 =\frac{|R\cap(A_{30}\setminus A_{20})|}{r_R}\geq0.
 \label{eq:formal-service-depth-monotonicity}
\end{equation}
Do đó các trung bình với cùng tập reference có nghĩa và cùng trọng số
không âm cũng không giảm. Không suy tăng Recall nghiêm ngặt, tăng
privacy hoặc ưu thế tại cùng bandwidth; tăng $L$ có giá phản hồi. Thay
$L_{\mathrm{plan}}$ làm Q khác đi thì lập luận giữ-cùng-Q không còn áp dụng.

\textbf{Sai vị trí local có thể phá tính đơn điệu, dù sai giống nhau ở
hai độ sâu.} Trên đường hai chiều, vị trí thật ở0; năm POI $a_1,\ldots,a_5$
ở1,2,3,4,5, còn $b$ ở10. Với $k=5$, reference thật là năm $a_i$.
Nếu pool nhỏ chứa chúng, local top5 trả cả năm dù ước lượng vị trí ở10.
Khi pool lớn thêm $b$, local ranking sai đưa $b$ lên đầu, đẩy một $a_i$
ra: Recall giảm từ1 xuống$4/5$. Đây là phản ví dụ về tập ứng viên lồng
nhau. Để có đúng prefix20/30, trên một graph có hướng có thể thêm15 POI
xa cả vị trí thật và vị trí ước lượng, nhưng nằm sau năm $a_i$ và trước
$b$ trong thứ tự từ Q. Chẳng hạn, ba trạng thái nguồn có các cạnh trực
tiếp tới POI với chi phí cho đúng ba thứ tự đó; các cạnh quay lại đủ dài
giữ những thứ tự và làm graph liên thông mạnh. Pool20 có năm $a_i$ và15
POI bổ sung; pool30 thêm $b$. Thu hồi thêm reference vẫn không giảm,
nhưng Recall \emph{sau xếp hạng sai} có thể giảm.

\textbf{Cache không tự phá tính đơn điệu.} Nếu hai cache bắt đầu giống
nhau, cùng lịch reset/expiry/version và cùng quy tắc hiệu lực, hợp các
phản hồi prefix giữ $A_{20}(t)\subseteq A_{30}(t)$ bằng quy nạp. Với
availability, trạng thái của ID chung còn phải nhất quán ở cùng epoch.
Ngược lại, so cache tích lũy L20 giữ POI $a$ từ quá khứ với L30 current-only
không còn $a$ thì các pool không lồng nhau. Eviction khác nhau cũng có
thể phá điều kiện này. Trạng thái cũ bị hết hạn trở thành unknown; không
được xem unknown là đang hoạt động hoặc là chắc chắn không hoạt động.

\textbf{Khác đồng hồ hoặc phiên bản không là so prefix.} Nếu phản hồi
L20 lấy trước cập nhật catalogue còn L30 lấy sau khi POI $a$ bị bỏ,
$a$ có thể nằm ở pool20 mà không có ở pool30. Reference hiện tại cũng
có thể đổi khi người dùng di chuyển hoặc availability đổi. Vì thế phải
giữ cùng thời điểm, miền, catalogue và chính sách cache khi kiểm tra
độ sâu; không lấy số liệu hai trạng thái dịch vụ khác nhau để suy một
bất đẳng thức từng event.

Các chứng minh trên giải thích điều kiện thiết kế và cách đọc phép đo.
Chúng không chứng minh posterior chính xác, GPS local không sai, dịch vụ
động có đủ trạng thái hiện hành hoặc mọi scenario đã được bảo vệ. Những
điểm đó được đánh giá riêng bằng benchmark và các diagnostic có phạm vi
nguồn, mẫu số và chi phí được ghi rõ.


\section{S9--S10, chức năng từng thành phần và giới hạn đóng góp}
\label{sec:current-endpoints-contribution}

Cấu hình hiện tại không warmup, không giữ truy vấn để gửi muộn và không bỏ phần cuối bằng holdback. Nó xử lý Q ngay tại các mốc đầu/cuối được khai báo. Không biết GPS tương lai để nhận diện ``60 s cuối''. Khi ứng dụng kết thúc, tín hiệu đóng phiên vẫn có thể quan sát được; mô hình bảo vệ tọa độ không tự che điều đó. Điểm đầu/cuối là thuộc tính của cả hành trình; (\ref{eq:current-transcript-bound}) không tự trở thành theorem riêng tư của nhà/nơi làm việc ẩn ngoài input đã đọc. Tấn công từ phần tuyến còn lại và siêu dữ liệu vẫn phải được đánh giá~\cite{dhondt2022endpoint}.

Thử nghiệm \texttt{GeoI-Endpoint20} tăng nhiễu ở các lần đọc được phép và gửi không delay là \emph{cấu hình riêng} đã có evidence endpoint. Nó không dùng detector biết trước điểm cuối, và không được tự gộp vào mô hình L30 hoặc chuyển số privacy từ cohort endpoint sang cohort mới. Việc xáo Q cũng không loại khả năng nối bằng Hungarian/ước lượng vận tốc. Chương thực nghiệm giữ các bank hữu hạn, protocol và cohort tách biệt.\footnote{Evidence endpoint không delay: \path{docs/research/2026-10-05_endpoint_noise.md}; tổng quát hóa và thứ tự: \path{docs/research/2026-10-05_endpoint_generalization.md}, \path{docs/research/2026-10-05_endpoint_order.md}; audit lý thuyết: \path{docs/research/2026-10-06_jisa_formal_audit.md}.}

\begin{table}[htbp]\centering\small
\caption{Vai trò theo kịch bản; ``hỗ trợ'' không là tuyên bố đã giải quyết hoàn toàn.}
\begin{tabularx}{\textwidth}{P{2.0cm}P{4.2cm}Y}\toprule
Kịch bản & Thành phần liên quan & Phạm vi thực tế\\\midrule
S1--S3 & REM, noisy reuse, bộ lọc; belief và miền đường & Cận tọa độ lý tưởng; chống suy luận từng điểm/chuỗi phải chấm attacker. Reuse chỉ giảm mẫu mới hữu hạn.\\
S4 & Sổ chung qua phiên & Chặn việc tự cấp lại cap trong epoch; không xóa account/IP hay tính nhận dạng quỹ đạo.\\
S5--S6 & Lịch sử đã bảo vệ, Q và ngân sách & Không nhận tương lai; attacker có thể dự đoán tuyến/đích từ hình học và routine.\\
S7 & Request mọi loại cố định; local ranking & Không thêm kênh purpose khi cùng $X$/lịch; intent tương quan và click vẫn là giới hạn.\\
S8 & Phân tích quan sát nhiều người & Chưa có cơ chế riêng chứng minh chống suy luận đồng hành.\\
S9--S10 & Nhiễu, cap, full-window attacker & Gửi không delay; không theorem endpoint hoặc che giờ đầu/cuối. Evidence Endpoint20 tách khỏi L30.\\\bottomrule
\end{tabularx}
\end{table}

Đóng góp là \textbf{ranh giới dữ liệu rõ ràng, tích hợp Geo-I với truy vấn có hướng và nhiều purpose cục bộ, cap qua phiên và đánh giá có kiểm tra độc lập}. Các primitive đã có tiền lệ; tăng $L$ chỉ điều chỉnh utility--cost, không chứng minh privacy tăng hoặc ưu thế tại cùng chi phí. Xác nhận synthetic cùng bản đồ chưa thay thế dữ liệu thực, chuyển đô thị, dịch vụ động, đo tài nguyên điện thoại hoặc chứng nhận sampler.
